\chapter{Proposed System}
\label{ch:design}

\section{Requirements}
\label{sec:design-requirements}

The system has two purposes. The first is to turn a public log of spaced-repetition reviews into a
leakage-free event table, on which several recall models can be fitted and scored in the same way. The second
is to run the proposed model's sampler reliably enough that its results can be checked. The requirements below
follow from these purposes. Each requirement is stated as something the system does, and each one is traceable
to a function or test in the code base.

Table~\ref{tab:functional} lists the functional requirements. They cover the pipeline from data download to the
report, and the checks that protect the evaluation from leakage.

\begin{table}[!htbp]
  \centering
  \caption{Functional requirements (FR).}
  \label{tab:functional}
  \footnotesize
  \begin{tabular}{@{}p{1.2cm}p{5.3cm}p{7.0cm}@{}}
    \toprule
    ID & Requirement & Realised by \\
    \midrule
    FR1 & Download the public review logs for a seeded sample of users & \texttt{download\_users(n\_users, seed, out\_dir)} \\
    FR2 & Load one user's log and apply the event rules & \texttt{load\_user(path)}: drop rating 0; keep state 4; drop cards whose first kept review is not new \\
    FR3 & Label each review with $y = 1[\text{rating} > 1]$ and recompute gaps from kept rows & \texttt{load\_user(path)}; \texttt{elapsed\_days}, \texttt{dt\_days} \\
    FR4 & Cap the number of cards per user with a seeded sample & \texttt{make\_split(df, max\_cards, seed)}; the reported run used \texttt{-{}-max\_cards 1000} (CLI default 300) \\
    FR5 & Split each user's reviews at a per-user time cutoff, and mark train and test events & \texttt{make\_split}: 80th percentile of \texttt{day\_offset} \\
    FR6 & Mark scored events, and keep same-day and first reviews as history only & \texttt{make\_split}: \texttt{scored} flag \\
    FR7 & Fit each model on training events only & \texttt{fit(events)} for FSRSModel, HLRModel, OursModel, LagLogitModel and, in Experiment~4, the v2, v3 and GBDT models \\
    FR8 & Predict one-step-ahead recall probabilities for scored events, with parameters frozen & \texttt{predict(events)} \\
    FR9 & Compute log loss, Brier score, AUC, ECE, RMSE over bins, mean bias and reliability data & \texttt{metrics.py} \\
    FR10 & Estimate user-clustered paired bootstrap intervals for metric differences & \texttt{paired\_bootstrap}: 1000 resamples, seed 0 \\
    FR11 & Save each model's predictions and a record of the run & \texttt{pred\_*.parquet}; \texttt{run\_\allowbreak info.json} \\
    FR12 & Resume an interrupted run without recomputing finished models & \texttt{-{}-resume}; saved predictions are reused only if their configuration matches \\
    FR13 & Write the report: metric tables, paired differences, user wins, plots & \texttt{report.py}: \texttt{metrics.csv}, \texttt{paired.csv}, \texttt{user\_wins.csv}, \texttt{report.md} \\
    FR14 & Report MCMC diagnostics for the sampler & Split $\hat{R}$, effective sample size, acceptance rates in \texttt{diagnostics\_*.json} \\
    FR15 & Show that flipping test labels does not change predictions & Leakage test in the test suite \\
    \bottomrule
  \end{tabular}
\end{table}

The functional requirements fall into three stages. FR1 to FR5 build the event table from the raw logs. FR6 to FR12
fit each model on the training events, predict the scored events and save the results. FR13 to FR15 produce the
report and check the two properties that the evaluation depends on, namely the sampler's diagnostics and the absence
of leakage.

Table~\ref{tab:nonfunctional} lists the non-functional requirements. Several of them are constraints that the
system does not fully meet. These are stated as they are, with the measured values. The sampler, in particular,
is slow and did not converge in the benchmark run.

\begin{table}[!htbp]
  \centering
  \footnotesize\setlength{\tabcolsep}{4pt}
  \caption{Non-functional requirements (NFR). Measured values are from the reported benchmark run.}
  \label{tab:nonfunctional}
  \footnotesize
  \begin{tabular}{@{}p{1.2cm}p{4.1cm}p{8.9cm}@{}}
    \toprule
    ID & Requirement & Status and evidence \\
    \midrule
    NFR1 & Reproducibility & Met for the given configuration. Card sampling, user sampling, the bootstrap and the sampler are all seeded (seed 0 in the reported run). The run configuration is recorded in \texttt{run\_\allowbreak info.json}. \\
    NFR2 & Correctness & Checked by \NumBenchTests{} automated pytest tests for the reported pipeline in \texttt{fsrs\_benchmark/\allowbreak tests}. The tests cover the split, the metrics against known values, the leakage check and the report. The \texttt{model\_diagnostics} package has 32 test functions. Four of its files need the optional packages arviz and polyagamma, which are not installed in the main environment, so those tests are not run by default. \\
    NFR3 & Evaluation cost & Met for the baselines, not for the proposed model. Fit times on the reported run: HLR 45 s, FSRS-6 optimised 653 s, proposed model 27,530 s (about 7.6 h on one core). Prediction for the proposed model takes 1,266 s. \\
    NFR4 & Memory and time per sampler sweep & Linear in the number of events: $O(N_{\text{events}})$ time and memory for $d = 1$. \\
    NFR5 & Hardware & A laptop with 8 cores and 19 GB RAM. CPU turbo was disabled for the long runs. \\
    NFR6 & Recoverability & Partly met. A finished model's predictions are kept, so an interrupted run resumes at the model level. The 8000-iteration sampler run was lost at sweep 1400 when the machine shut down, and it has no result. \\
    NFR7 & Convergence of the proposed model's sampler & Not met. The gate was a split $\hat{R}$ below 1.01 and an effective sample size (ESS) of at least 400. In the learned-rate run, seven scalars fail the gate: $\lambda$ ($\hat{R}$ 2.44, ESS 3.0), $\sigma^2$ ($\hat{R}$ 2.07, ESS 3.3), $b_0$ ($\hat{R}$ 1.85, ESS 3.9), $\tau^2_\beta$ ($\hat{R}$ 1.88, ESS 4.8), $g_0$ (ESS 80.6), $\tau^2_\alpha$ (ESS 250) and $\tau^2_\gamma$ ($\hat{R}$ 1.06, ESS 143). The $+\kappa$ run also fails, including $\kappa$ ($\hat{R}$ 1.94, ESS 3.5). The scores are therefore provisional. \\
    NFR8 & Transparency & Met. Every metric is written from scratch and tested against known values. \\
    \bottomrule
  \end{tabular}
\end{table}

\section{System Architecture}
\label{sec:design-architecture}

The system is organised in four layers, as shown in Figure~\ref{fig:tk_architecture}. Each layer reads
the output of the layer before it and writes to a store on disk, so that a later layer can be rerun without
repeating the earlier ones.

\begin{figure}[htbp]
  \centering
  % tk_architecture: four layers and the stores they read and write (DIAGRAMS.md item 1)
\begin{tikzpicture}[font=\small]
  % layer 1: data ingestion (box y0=-2.6)
  \draw[rounded corners=4pt, draw=black!60, fill=black!2] (0,-2.6) rectangle (10.8,0);
  \node[anchor=west, font=\small\bfseries] at (0.25,-0.35) {1. Data ingestion};
  \node[box, text width=2.8cm, minimum height=1.3cm] at (1.8,-1.6) {HF download\\\texttt{download\_\allowbreak users}};
  \node[box, text width=2.8cm, minimum height=1.3cm] at (5.4,-1.6) {\texttt{load\_user}};
  \node[box, text width=2.8cm, minimum height=1.3cm] at (9.0,-1.6) {Event rules\\(filter, label $y$)};
  % layer 2: split and protocol (y0=-6.1)
  \draw[rounded corners=4pt, draw=black!60, fill=black!2] (0,-6.1) rectangle (10.8,-3.5);
  \node[anchor=west, font=\small\bfseries] at (0.25,-3.85) {2. Split and protocol};
  \node[box, text width=2.1cm, minimum height=1.3cm] at (4.05,-5.1) {Per-user 80th-pct\\cutoff};
  \node[box, text width=2.1cm, minimum height=1.3cm] at (1.45,-5.1) {Card cap\\($\le 1000$ per user)};
  \node[box, text width=2.1cm, minimum height=1.3cm] at (6.65,-5.1) {Scored-event\\filter};
  \node[box, text width=2.1cm, minimum height=1.3cm] at (9.25,-5.1) {One-step-ahead\\protocol};
  % layer 3: models (y0=-9.6), coloured per PLAN.md
  \draw[rounded corners=4pt, draw=black!60, fill=black!2] (0,-9.6) rectangle (10.8,-7.0);
  \node[anchor=west, font=\small\bfseries] at (0.25,-7.35) {3. Models};
  \node[box, modelbox=fsrsdef, text width=1.45cm, inner sep=1.5pt, font=\footnotesize, minimum height=1.3cm] at (1.1,-8.6) {\mbox{FSRS-6}\\default};
  \node[box, modelbox=fsrsopt, text width=1.45cm, inner sep=1.5pt, font=\footnotesize, minimum height=1.3cm] at (2.85,-8.6) {\mbox{FSRS-6}\\optimised};
  \node[box, modelbox=hlr, text width=1.45cm, inner sep=1.5pt, font=\footnotesize, minimum height=1.3cm] at (4.6,-8.6) {HLR};
  \node[box, modelbox=ours, text width=1.45cm, inner sep=1.5pt, font=\footnotesize, minimum height=1.3cm] at (6.35,-8.6) {Ours\\(learned\\$\lambda$)};
  \node[box, modelbox=ourskappa, text width=1.45cm, inner sep=1.5pt, font=\footnotesize, minimum height=1.3cm] at (8.1,-8.6) {Ours\\(+$\kappa$)};
  \node[box, modelbox=laglogit, text width=1.45cm, inner sep=1.5pt, font=\footnotesize, minimum height=1.3cm] at (9.85,-8.6) {LagLogit\\(Exp.~3)};
  % layer 4: evaluation (y0=-13.1)
  \draw[rounded corners=4pt, draw=black!60, fill=black!2] (0,-13.1) rectangle (10.8,-10.5);
  \node[anchor=west, font=\small\bfseries] at (0.25,-10.85) {4. Evaluation};
  \node[box, text width=2.8cm, minimum height=1.3cm] at (1.8,-12.1) {Metrics\\(log loss, Brier,\\AUC, ECE)};
  \node[box, text width=2.8cm, minimum height=1.3cm] at (5.4,-12.1) {Paired bootstrap\\(user-clustered)};
  \node[box, text width=2.8cm, minimum height=1.3cm] at (9.0,-12.1) {Report\\(\texttt{metrics.csv}, plots)};
  % arrows between layers
  \draw[arrow] (5.4,-2.6) -- (5.4,-3.5);
  \draw[arrow] (5.4,-6.1) -- (5.4,-7.0);
  \draw[arrow] (5.4,-9.6) -- (5.4,-10.5);
  % side stores
  \node[store, minimum width=2.6cm, minimum height=1.3cm] (hf) at (12.9,-1.3) {anki-revlogs-10k\\(HF)};
  \node[store, minimum width=2.6cm, minimum height=1.3cm] (ev) at (12.9,-4.8) {events.parquet};
  \node[store, minimum width=2.6cm, minimum height=1.3cm] (pr) at (12.9,-8.3) {pred\_*.parquet};
  \draw[arrow] (hf.west) -- (10.8,-1.3);
  \draw[arrow] (10.8,-4.8) -- (ev.west);
  \draw[arrow] (10.8,-8.3) -- (pr.west);
\end{tikzpicture}

  \caption{Four layers of the benchmark system, with the stores they read and write. Layer 1 downloads and
    filters the review logs. Layer 2 builds the leakage-free split. Layer 3 fits and predicts with the models of Experiments 2 and 3 (LagLogit is used on Duolingo only), shown in their plot colours; the Experiment 4 models are added to this layer in Chapter~\ref{ch:results}. Layer 4 computes the metrics and the report.}
  \label{fig:tk_architecture}
\end{figure}

Layer 1 is the data ingestion layer. It downloads the anki-revlogs-10k dataset from Hugging Face for the Anki benchmark, and it reads the public Duolingo trace file for Experiment~3. It loads each
user's review file, and applies the event rules. Its output is one table of kept reviews per user. Layer 2
applies the card cap, computes each user's cutoff, and writes the event table to \texttt{events.parquet}.
Layer 3 holds the five models of Experiments 2 and 3: FSRS-6 with default parameters, FSRS-6 with optimised parameters, HLR, the
proposed model with a learned rate, and the proposed model with the practice term $\kappa$ (Duolingo adds a lag-logistic baseline). Each model is fitted
on the training events and writes its one-step-ahead predictions to a \texttt{pred\_*.parquet} file. Layer 4
reads the predictions together with the event table, computes the metrics, and produces the paired bootstrap
intervals and the report. Layer 3 is extended in Experiment~4 by three further models, \texttt{ours\_v2}, \texttt{ours\_v3} and \texttt{gbdt}, which use the same \texttt{fit} and \texttt{predict} interface and write the same prediction files.

The architecture keeps the models separate from the evaluation. A model is any object with the same two
methods, \texttt{fit} and \texttt{predict}. The evaluation layer never inspects a model's internals, so a new
model can be added by writing one class and one key.

\section{Data-Flow Diagrams}
\label{sec:design-dfd}

Figure~\ref{fig:tk_dfd0} is the context diagram. It shows the benchmark system as a single process, with the
researcher on the left, the Hugging Face dataset above, and a future spaced-repetition application on the right.
The researcher supplies the configuration as command-line arguments and receives the predictions and report. The
dataset supplies the review logs. The application is dashed because it is future work, and the recall
probabilities flowing to it are not implemented.

\begin{figure}[htbp]
  \centering
  % tk_dfd0: context diagram (DIAGRAMS.md item 2)
\begin{tikzpicture}[font=\small]
  \node[process, minimum size=3.8cm, font=\small] (sys) at (7.0,0) {0. Retention\\Benchmark System};
  \node[entity, minimum width=3.0cm, minimum height=1.2cm] (res) at (1.6,0) {Researcher};
  \node[entity, minimum width=3.6cm, minimum height=1.0cm] (hf) at (7.0,4.2) {Hugging Face\\dataset};
  \node[entity, minimum width=2.8cm, minimum height=1.2cm, dashed] (app) at (12.7,0) {Spaced-repetition\\app (future)};
  % config in, report out
  \draw[arrow] (3.1,0.4) -- (5.14,0.4) node[above, midway, align=center, font=\footnotesize] {config /\\CLI args};
  \draw[arrow] (5.14,-0.4) -- (3.1,-0.4) node[below, midway, align=center, font=\footnotesize] {predictions\\\& report};
  % revlogs in
  \draw[arrow] (7.0,3.7) -- (7.0,1.9) node[right, midway, font=\footnotesize, xshift=3pt] {revlogs};
  % recall probabilities out (future app, dashed)
  \draw[dashedarrow] (8.86,0) -- (11.0,0) node[above, midway, align=center, font=\footnotesize] {recall\\probabilities};
\end{tikzpicture}

  \caption{Context diagram (level 0). The benchmark system exchanges the configuration, the review logs, the
    predictions and the report with its environment. The dashed application is future work.}
  \label{fig:tk_dfd0}
\end{figure}

Figure~\ref{fig:tk_dfd1} decomposes the system into five processes and four data stores, in the
Yourdon--DeMarco notation. Processes are circles, data stores are open-ended rectangles, and external entities are
rectangles. Process 1.0 downloads and cleans the raw revlogs into store D1. Process 2.0 builds the split, and
writes the event table to store D2. Process 3.0 fits the models on the training events and writes their fitted
parameters or posterior draws to store D3. Process 4.0 predicts the scored events, using the event table and the
fitted parameters, and writes the predictions to store D4. Process 5.0 evaluates the predictions against the labels and user identifiers from D2, and writes the report for the researcher.

\begin{figure}[htbp]
  \centering
  % tk_dfd1: level-1 data flow diagram, Yourdon/DeMarco notation (DIAGRAMS.md item 3)
\begin{tikzpicture}[font=\small]
  % column x: c1=1.3, c2=4.2, c3=7.0, c4=9.8, c5=12.6; row y: A=0, B=-3.4, C=-6.8, D=-10.2
  % row A
  \node[entity, minimum width=2.4cm, minimum height=1.0cm] (hf) at (1.3,0) {Hugging Face\\dataset};
  \node[process, minimum size=2.2cm] (p1) at (4.2,0) {1.0\\Download \&\\clean};
  \node[dstore, minimum width=2.4cm, minimum height=1.0cm] (d1) at (7.0,0) {D1\\raw revlogs};
  \node[process, minimum size=2.2cm] (p2) at (9.8,0) {2.0\\Split};
  % row B
  \node[dstore, minimum width=2.4cm, minimum height=1.0cm] (d2) at (9.8,-3.4) {D2\\events.parquet};
  \node[process, minimum size=2.2cm] (p4) at (7.0,-3.4) {4.0\\Predict\\scored events};
  \node[dstore, minimum width=2.4cm, minimum height=1.0cm] (d4) at (4.2,-3.4) {D4\\pred\_*.parquet};
  \node[process, minimum size=2.2cm] (p3) at (12.6,-3.4) {3.0\\Fit models\\(train)};
  % row C
  \node[dstore, minimum width=2.8cm, minimum height=1.0cm] (d3) at (12.6,-6.8) {D3\\fitted params /\\posterior draws};
  \node[process, minimum size=2.2cm] (p5) at (4.2,-6.8) {5.0\\Evaluate \&\\report};
  % row D
  \node[entity, minimum width=2.4cm, minimum height=1.0cm] (res) at (4.2,-10.2) {Researcher};

  % row A flows
  \draw[arrow] (hf.east) -- (p1.west);
  \draw[arrow] (p1.east) -- (d1.west);
  \draw[arrow] (d1.east) -- (p2.west);
  % 2.0 -> D2 (down), D2 -> 3.0 (right), D2 -> 4.0 (left)
  \draw[arrow] (p2.south) -- (d2.north) node[right, midway, font=\footnotesize, xshift=3pt] {event table};
  \draw[arrow] (d2.east) -- (p3.west);
  \draw[arrow] (d2.west) -- (p4.east);
  % 3.0 -> D3 (down), D3 -> 4.0 (left, then up)
  \draw[arrow] (p3.south) -- (d3.north) node[right, midway, font=\footnotesize, xshift=3pt, align=left] {fitted\\params};
  \draw[arrow] (d3.west) -- (7.0,-6.8) -- (p4.south);
  \node[font=\footnotesize, anchor=south] at (8.4,-6.75) {posterior draws};
  % 4.0 -> D4 (left), D4 -> 5.0 (down)
  \draw[arrow] (p4.west) -- (d4.east);
  \draw[arrow] (d4.south) -- (p5.north) node[right, midway, font=\footnotesize, xshift=3pt] {predictions};
  % 5.0 -> Researcher, with the output files
  % D2 -> 5.0: process 5.0 reads the labels and user ids from the event table, routed below the D3 path
  \draw[arrow] (d2.south) -- (9.8,-8.15) -- (5.6,-8.15) -- (p5.-45);
  \node[font=\footnotesize, align=center] at (7.9,-7.85) {labels, user ids};
  \draw[arrow] (p5.south) -- (res.north) node[right, midway, align=left, font=\footnotesize, xshift=4pt] {report.md,\\metrics.csv};

  % open-ended stores: draw only the two horizontal edges
  \foreach \s in {d1,d2,d3,d4} {
    \draw[black!70, line width=0.5pt] (\s.north west) -- (\s.north east) (\s.south west) -- (\s.south east);
  }
\end{tikzpicture}

  \caption{Level-1 data-flow diagram. Processes 1.0 to 5.0 correspond to the four layers of
    Figure~\ref{fig:tk_architecture}. Data stores D1 to D4 hold the raw logs, the event table, the fitted
    parameters and the predictions.}
  \label{fig:tk_dfd1}
\end{figure}

Reading the diagram from the top left, the flow is a pipeline with one branch from process 3.0. The fitted parameters
from process 3.0 are needed in process 4.0, which is why store D3 feeds back into the prediction step. The
process 5.0 reads the predictions in D4 and the labels and user identifiers from D2. The evaluation therefore depends on the event table and the predictions, and not on the models themselves. This is the property that keeps the metrics comparable.

\section{Unified Modelling Language Models}
\label{sec:design-uml}

\subsection{Use cases}

Figure~\ref{fig:tk_usecase} shows the use cases. The system boundary contains eight use cases, from downloading
the data to generating the report. The researcher is the primary actor and is associated with every use case.
The dataset provider is associated with the download use case. A learner application is shown dashed, because it
would consume the predicted recall probabilities in a future version. The use case \emph{Resume interrupted run}
extends \emph{Fit model}, because it is an optional path that skips a model whose predictions already exist.

\begin{figure}[htbp]
  \centering
  % tk_usecase: UML use-case diagram (DIAGRAMS.md item 4)
% All use cases sit in one column so the Researcher can be associated with each
% of them without any association line crossing another ellipse.
\begin{tikzpicture}[font=\small]
  % system boundary
  \draw[black!60, line width=0.6pt] (1.6,0.2) rectangle (8.8,9.7);
  \node[font=\small\bfseries] at (5.2,9.3) {Benchmark \& Retention Model};
  \node[usecase, minimum width=4.0cm] (uc1) at (4.6,8.6) {Download data};
  \node[usecase, minimum width=4.0cm] (uc2) at (4.6,7.5) {Build event table \& split};
  \node[usecase, minimum width=4.0cm] (uc3) at (4.6,6.4) {Fit model};
  \node[usecase, minimum width=4.0cm] (uc4) at (4.6,5.3) {Resume interrupted run};
  \node[usecase, minimum width=4.0cm] (uc5) at (4.6,4.2) {Run MCMC diagnostics};
  \node[usecase, minimum width=4.0cm] (uc6) at (4.6,3.1) {Predict recall probability};
  \node[usecase, minimum width=4.0cm] (uc7) at (4.6,2.0) {Evaluate \& compare models};
  \node[usecase, minimum width=4.0cm] (uc8) at (4.6,0.9) {Generate report};
  % actors
  \StickActor{(1.0,4.6)}{}
  % label to the left of the figure, clear of every association line
  \node[actor, anchor=east] at (0.75,4.95) {Researcher};
  \StickActor{(10.2,8.6)}{Dataset Provider}
  \StickActor{(10.2,3.1)}{Learner App\\(future)}
  % Researcher associations: one fan from the actor to each use case
  \foreach \u in {uc1,uc2,uc3,uc4,uc5,uc6,uc7,uc8} {
    \draw[black!75, line width=0.5pt] (1.45,4.6) -- (\u.west);
  }
  % Dataset provider supplies the revlogs
  \draw[black!75, line width=0.5pt] (9.6,8.6) -- (uc1.east);
  % learner app consumes the predictions (future)
  \draw[black!75, line width=0.5pt, dashed] (9.6,3.1) -- (uc6.east);
  % <<extend>>: Resume interrupted run extends Fit model
  \draw[dashed, line width=0.5pt, -{Triangle[open, length=2.6mm, width=2.0mm]}]
    (uc4.east) -- ++(0.7,0) |- node[pos=0.25, above right, font=\footnotesize] {\guillemotleft extend\guillemotright} (uc3.east);
\end{tikzpicture}

  \caption{Use-case diagram. The researcher uses all eight use cases. The dataset provider supplies the review logs.
    The learner application is future work. \emph{Resume interrupted run} extends \emph{Fit model}.}
  \label{fig:tk_usecase}
\end{figure}

\subsection{Classes}

Figure~\ref{fig:tk_class} shows the main classes. The interface \emph{Model} declares the two methods, \texttt{fit}
and \texttt{predict}. Four classes realise it: \texttt{FSRSModel}, which has the options \texttt{optimize} and
\texttt{n\_workers}; \texttt{HLRModel}, with the learning rate and the L2 weight; \texttt{OursModel}, with the
options that control the sampler, namely the switch \texttt{use\_kappa} for the $+\kappa$ variant, the number of iterations, the burn-in, the thinning, the seed and the number of posterior draws used for prediction; and \texttt{LagLogitModel}, which has one option that selects the
time term. The realisation arrows are dashed, as UML requires.

\texttt{OursModel} is a composition of \texttt{DynamicRetention\allowbreak{}ModelDT}, which holds the latent state, the
parameters and the sampler steps. The composition is drawn with a filled diamond. The diagram shows five of the class's methods: \texttt{run}, \texttt{\_ffbs}, \texttt{\_update\_effects}, \texttt{\_update\_lam\allowbreak\_sigma2} and \texttt{predict}. This class uses the P\'olya--Gamma sampler \texttt{sample\allowbreak\_pg1}
and the Numba kernel \texttt{\allowbreak\_ffbs\allowbreak\_step\allowbreak\_numba}. These dependencies are drawn as dashed arrows, because they are
calls rather than owned parts.

\begin{figure}[htbp]
  \centering
  % tk_class: UML class diagram of the model classes (DIAGRAMS.md item 5)
\begin{tikzpicture}[font=\footnotesize]
  % interface
  \node[uml, rectangle split, rectangle split parts=2, text width=4.0cm, anchor=north,
        align=left, inner sep=3pt] (iface) at (7.125,10.6) {
    \centering\guillemotleft interface\guillemotright\\\centering\textbf{\textit{Model}}
    \nodepart{two}
    \texttt{fit(events)}\\ \texttt{predict(events)}
  };
  % implementing classes (top at y=8.2)
  \node[uml, rectangle split, rectangle split parts=3, text width=3.3cm, anchor=north, inner sep=3pt]
    (fsrs) at (1.65,7.5) {
    \centering\textbf{FSRSModel}
    \nodepart{two}
    \texttt{optimize}\\ \texttt{n\_workers}
    \nodepart{three}
    \texttt{fit(events)}\\ \texttt{predict(events)}
  };
  \node[uml, rectangle split, rectangle split parts=3, text width=3.3cm, anchor=north, inner sep=3pt]
    (hlr) at (5.3,7.5) {
    \centering\textbf{HLRModel}
    \nodepart{two}
    \texttt{lrate}\\ \texttt{l2wt}
    \nodepart{three}
    \texttt{fit(events)}\\ \texttt{predict(events)}
  };
  \node[uml, rectangle split, rectangle split parts=3, text width=3.3cm, anchor=north, inner sep=3pt]
    (ours) at (8.95,7.5) {
    \centering\textbf{OursModel}
    \nodepart{two}
    \texttt{use\_kappa}\\ \texttt{iters}\\ \texttt{burn}\\ \texttt{thin}\\ \texttt{seed}\\ \texttt{max\_draws}
    \nodepart{three}
    \texttt{fit(events)}\\ \texttt{predict(events)}
  };
  \node[uml, rectangle split, rectangle split parts=3, text width=3.3cm, anchor=north, inner sep=3pt]
    (lag) at (12.6,7.5) {
    \centering\textbf{LagLogitModel}
    \nodepart{two}
    \texttt{use\_dt}
    \nodepart{three}
    \texttt{fit(events)}\\ \texttt{predict(events)}
  };
  % realisation: one bus per class to the interface
  \draw[realize] (fsrs.north) -- ++(0,0.9) -| (iface.south);
  \draw[realize] (hlr.north) -- ++(0,0.9) -| (iface.south);
  \draw[realize] (ours.north) -- ++(0,0.9) -| (iface.south);
  \draw[realize] (lag.north) -- ++(0,0.9) -| (iface.south);
  % core model (composition from OursModel)
  \node[uml, rectangle split, rectangle split parts=3, text width=8.0cm, anchor=north, inner sep=3pt]
    (dr) at (6.2,2.6) {
    \centering\textbf{DynamicRetentionModelDT}
    \nodepart{two}
    \texttt{lam, sigma2, kappa, b0, alpha, beta, gamma, tau2}
    \nodepart{three}
    \texttt{run()}\\ \texttt{\_ffbs()}\\ \texttt{\_update\_effects()}\\ \texttt{\_update\_lam\_sigma2()}\\ \texttt{predict()}
  };
  \draw[composition] (ours.south) -- ++(0,-0.5) -| (dr.north);
  % helpers used by the core model
  \node[uml, text width=3.6cm, align=center, anchor=north, inner sep=3pt] (pg) at (3.2,-1.8) {
    \guillemotleft function\guillemotright\\\texttt{sample\_pg1}\\(pg\_sampler)
  };
  \node[uml, text width=4.2cm, align=center, anchor=north, inner sep=3pt] (ffbs) at (10.4,-1.8) {
    \guillemotleft kernel\guillemotright\\\texttt{\_ffbs\_step\_numba}\\(numba)
  };
  \draw[dashedarrow] (pg.north |- dr.south) -- (pg.north) node[right, midway, font=\footnotesize, xshift=3pt] {\guillemotleft uses\guillemotright};
  \draw[dashedarrow] (ffbs.north |- dr.south) -- (ffbs.north) node[right, midway, font=\footnotesize, xshift=3pt] {\guillemotleft uses\guillemotright};
\end{tikzpicture}

  \caption[Class diagram of the model classes]{Class diagram of the models of Experiments 2 and 3 (the Experiment~4 models share the same interface and are not drawn). Four model classes realise the \emph{Model} interface. \texttt{OursModel} is a composition
    of \texttt{DynamicRetention\allowbreak ModelDT}, which uses the P\'olya--Gamma sampler and FFBS.}
  \label{fig:tk_class}
\end{figure}

\subsection{Run-time sequence}

Figure~\ref{fig:tk_sequence} shows the run-time sequence of \texttt{run.py}. The researcher calls \texttt{run} with the
arguments. The run loads the event source and builds the event table through the split module. It then loops over the
model keys. For each key it creates the model, fits it on the training events, predicts the scored events, and saves
the predictions. The loop is skipped for a key whose file already exists when \texttt{-{}-resume} is given. Finally the
run builds the report, which writes \texttt{metrics.csv}, \texttt{paired.csv} and \texttt{report.md}.

\begin{figure}[htbp]
  \centering
  % tk_sequence: run-time sequence of run.py (DIAGRAMS.md item 6), pgf-umlsd
\begin{sequencediagram}
  % message labels sit on a white background, so no lifeline runs through them
  \tikzset{every node/.append style={fill=white, inner sep=1pt}}
  \newinst[0.75]{res}{Researcher}
  \newinst[0.75]{run}{\texttt{run.py}}
  \newinst[0.75]{data}{\texttt{data/split}}
  \newinst[0.75]{mod}{Model}
  \newinst[0.75]{rep}{\texttt{metrics/report}}

  \begin{call}{res}{run(args)}{run}{}\end{call}
  \begin{call}{run}{load events}{data}{\raisebox{2.5pt}{events}}\end{call}
  \begin{call}{run}{make split}{data}{\raisebox{2.5pt}{event table}}\end{call}

  \begin{sdblock}{loop}{\hspace{4pt}\footnotesize for each model key (skip if -{}-resume)}
    \begin{call}{run}{make model}{mod}{\raisebox{2.5pt}{model}}\end{call}
    \begin{call}{run}{fit(train)}{mod}{}\end{call}
    \begin{call}{run}{predict(scored)}{mod}{\raisebox{2.5pt}{probabilities}}\end{call}
    \begin{callself}[2]{run}{save pred\_<key>.parquet}{}\end{callself}
  \end{sdblock}

  \begin{call}{run}{build report}{rep}{\raisebox{2.5pt}{metrics.csv, paired.csv, report.md}}\end{call}
\end{sequencediagram}

  \caption{Sequence diagram of \texttt{run.py}. The loop over models is the only repeated part. The optional
    \emph{skip if -{}-resume} guard avoids refitting a model whose predictions exist.}
  \label{fig:tk_sequence}
\end{figure}

\subsection{Gibbs sweep}

Figure~\ref{fig:tk_activity_sweep} is the activity diagram of one sweep of the sampler. The order of the steps is the
order in which \texttt{run()} executes them. The sweep first computes the transition coefficients $\phi$ and $q$ for
each gap. It then draws the P\'olya--Gamma variables, the latent state by FFBS, the Gaussian effects, the learner and
card variances, and the rate $\lambda$ by a Metropolis--Hastings step on $\log\lambda$. If the practice term is
used, a Metropolis--Hastings step on $\kappa$ follows. The sweep ends with the draw of $\sigma^2$. A draw is stored
only after burn-in and only on a thinning step. The loop continues until the requested number of iterations is
reached.

\begin{figure}[htbp]
  \centering
  % tk_activity_sweep: one Gibbs sweep, in the order of FACTS.md section 1 (DIAGRAMS.md item 7)
\begin{tikzpicture}[font=\footnotesize]
  \tikzset{actstep/.style={box, text width=5.6cm, minimum height=0.7cm, font=\footnotesize}}
  \node[circle, fill=black, inner sep=0pt, minimum size=3.5mm] (start) at (0,0) {};
  \node[actstep] (a1) at (0,-1.2) {Compute $\phi$, $q$ (\texttt{\_dyn})};
  \node[actstep] (a2) at (0,-2.3) {Draw $\omega \sim \mathrm{PG}(1,\psi)$};
  \node[actstep] (a3) at (0,-3.4) {FFBS: draw $\psi \mid \omega, \theta$};
  \node[actstep] (a4) at (0,-4.5) {Gaussian draws: $b_0, \alpha, \beta, \gamma, g_0$};
  \node[actstep] (a5) at (0,-5.6) {Inverse-Gamma draws: $\tau^2_\alpha, \tau^2_\beta, \tau^2_\gamma$};
  \node[actstep] (a6) at (0,-6.7) {Metropolis--Hastings on $\log\lambda$ (5 proposals)};
  \node[decision, minimum width=2.8cm, minimum height=0.9cm] (d1) at (0,-7.9) {use $\kappa$?};
  \node[actstep, text width=3.2cm] (a7) at (5.2,-7.9) {MH on $\kappa$ (5 proposals)};
  \node[actstep, text width=5.6cm] (a8) at (0,-9.2) {Inverse-Gamma draw $\sigma^2$};
  \node[decision, minimum width=6.4cm, minimum height=1.2cm] (d2) at (0,-10.7) {it $\ge$ burn-in and\\on thin step?};
  \node[actstep, text width=3.2cm] (a9) at (6.4,-10.7) {Store draw};
  \node[decision, minimum width=4.2cm, minimum height=0.9cm] (d3) at (0,-12.7) {more iterations?};
  \node[circle, draw=black, line width=0.6pt, minimum size=4mm, inner sep=0pt] (end) at (0,-14.2) {};
  \node[circle, fill=black, minimum size=2mm, inner sep=0pt] at (0,-14.2) {};

  % main flow
  \draw[arrow] (start) -- (a1);
  \draw[arrow] (a1) -- (a2);
  \draw[arrow] (a2) -- (a3);
  \draw[arrow] (a3) -- (a4);
  \draw[arrow] (a4) -- (a5);
  \draw[arrow] (a5) -- (a6);
  \draw[arrow] (a6) -- (d1);
  % use kappa? yes -> MH on kappa -> merge into sigma^2 step; no -> straight down
  \draw[arrow] (d1.east) -- (a7.west) node[above, midway, font=\footnotesize] {yes};
  \draw[arrow] (a7.south) |- (a8.east);
  \draw[arrow] (d1.south) -- node[right, font=\footnotesize, xshift=2pt] {no} (a8.north);
  \draw[arrow] (a8) -- (d2);
  % burn-in and thin: yes -> store draw; both paths reach "more iterations?"
  \draw[arrow] (d2.east) -- (a9.west) node[above, midway, font=\footnotesize] {yes};
  \draw[arrow] (a9.south) |- (d3.east);
  \draw[arrow] (d2.south) -- node[right, font=\footnotesize, xshift=2pt] {no} (d3.north);
  % more iterations? yes -> back to the top of the sweep; no -> end
  \draw[arrow] (d3.south) -- node[right, font=\footnotesize, xshift=2pt] {no} (end);
  \draw[arrow] (d3.west) -- ++(-2.0,0) node[above, midway, font=\footnotesize] {yes} |- (a1.west);
\end{tikzpicture}

  \caption{Activity diagram of one Gibbs sweep of the proposed model. The order matches the implementation. The
    $\kappa$ step is taken only in the variant with the practice term. Draws are stored after burn-in on thinning
    steps.}
  \label{fig:tk_activity_sweep}
\end{figure}
